% 每行仅保留用于区分接入位置的模块，不展开完整Transformer层。
\begin{tikzpicture}[x=1cm,y=1cm,font=\figurefont\small,
  block/.style={draw=FigureModel,fill=FigureModel!10,line width=0.8pt,
    minimum height=9mm,minimum width=21mm,align=center},
  input/.style={block,draw=FigureInput,fill=FigureInput!10},
  output/.style={block,draw=FigureOutput,fill=FigureOutput!10},
  arr/.style={-{Stealth[length=2mm]},BookInk,line width=0.8pt}]
  \node[anchor=west,font=\figurefont\small\bfseries] at (0,6.7) {BLIP-2：查询压缩后作为输入前缀};
  \node[input] (bf) at (1.1,5.7) {视觉特征};
  \node[block] (bq) at (4.1,5.7) {Q-Former};
  \node[block] (bp) at (7.2,5.7) {投影后的\\视觉前缀};
  \node[block] (bl) at (10.7,5.7) {语言模型};
  \node[output] (bo) at (14.1,5.7) {回答词元};
  \draw[arr] (bf) -- (bq);
  \draw[arr] (bq) -- (bp);
  \draw[arr] (bp) -- (bl);
  \draw[arr] (bl) -- (bo);
  \node[anchor=north] at (4.1,5.15) {可学习查询读取图像};
  \node (btext) at (10.7,4.7) {文字问题};
  \draw[arr] (btext) -- (bl);

  \node[anchor=west,font=\figurefont\small\bfseries] at (0,3.9) {Flamingo：在语言层之间读取视觉表示};
  \node[input] (ff) at (1.1,2.45) {视觉特征};
  \node[block] (fr) at (4.1,2.45) {Perceiver\\Resampler};
  \node[block] (fc) at (7.2,2.45) {门控交叉\\注意力};
  \node[block] (fl) at (10.7,2.45) {后续语言层};
  \node[output] (fo) at (14.1,2.45) {回答词元};
  \draw[arr] (ff) -- (fr);
  \draw[arr] (fr) -- node[above] {$K,V$} (fc);
  \draw[arr] (fc) -- (fl);
  \draw[arr] (fl) -- (fo);
  \node (fh) at (7.2,3.35) {语言隐状态$Q$};
  \draw[arr] (fh) -- (fc);
  \node[anchor=north] at (4.1,1.9) {每个图像／视频64个单元};

  \node[anchor=west,font=\figurefont\small\bfseries] at (0,0.7) {原始LLaVA：网格特征投影后进入共享序列};
  \node[input] (lf) at (1.1,-0.25) {视觉网格};
  \node[block] (lp) at (4.1,-0.25) {线性投影};
  \node[block] (ls) at (7.2,-0.25) {视觉单元\\与问题序列};
  \node[block] (ll) at (10.7,-0.25) {语言模型};
  \node[output] (lo) at (14.1,-0.25) {回答词元};
  \draw[arr] (lf) -- (lp);
  \draw[arr] (lp) -- (ls);
  \draw[arr] (ls) -- (ll);
  \draw[arr] (ll) -- (lo);
\end{tikzpicture}
